\section{Shared Data Engine and Foundation Pretraining}
\label{sec:data}

\subsection{A Common Data Production Pipeline}

StepAudio 2.5 adopts an automated data production pipeline that jointly supports speech understanding, TTS, and speech dialogue tasks. Raw audio is first processed with sound event detection (SED) and voice activity detection (VAD) to filter low-quality non-speech segments. Adjacent VAD segments are then merged and re-segmented into base samples with relatively complete semantics and suitable duration. For each audio clip within the base samples, audio-level annotations are performed, including audio quality scoring, synthetic voice detection, and speaker count labeling. At the text annotation level, dual ASR models are employed for transcription and language identification. The resulting transcripts are cross-validated with metrics such as WER, edit distance, and speech rate. Based on the ASR transcription, semantic completeness assessment and content classification are further carried out for each base sample. Finally, according to metadata, the data is categorized and graded by language, duration, semantic quality score, and audio quality score, enabling the pretraining phase to sample different data qualities for different training stages.

\subsection{Progressive Foundation Training}

StepAudio 2.5 is initialized from a textual MoE LLM and then continually pre-trained on 2.2T tokens of text and audio data. The training curriculum follows a concrete staged recipe rather than a loosely defined scaling process.

The first stage follows Step-Audio 2 and uses 3B tokens of ASR data to align speech and text feature spaces within the adaptor. During this alignment phase, both the audio encoder and the LLM remain frozen, and only the adaptor is trained. This stage establishes the initial interface through which acoustic features can be consumed by the text-native decoder.

After alignment, the model vocabulary is expanded with speech tokens, and unified multimodal training begins with a sequence length of 16K. This main pretraining mixture contains 800B tokens of text data and 800B tokens of speech data. The speech portion includes ASR, TTS, speech-to-text translation, utterance-level text-speech interleaved continuation, and speech-to-speech conversation data. In other words, the model is not exposed to audio only as transcription input, but as a general sequence modality appearing in multiple input-output configurations.

This multimodal phase is itself divided into two stages. The first is a 128B-token warmup stage designed to stabilize the newly introduced speech vocabulary and help the MoE experts adapt quickly to audio-modality data. In this stage, the adaptor, embedding layer, and output layer use larger learning rates than the base model, while the MoE router uses a smaller learning rate to reduce disruption to the text modality. The second is the main training stage, where these layer-specific learning rates are brought back in line with the base learning rate, and the MoE auxiliary loss coefficient together with the router learning rate are progressively annealed to maintain a better balance between expert utilization and top-$k$ routing probabilities.

Finally, the model enters a cooldown phase on 600B tokens of high-quality text and audio data, with the sequence length increased to 32K. In addition to the data types already used in the main training stage, this phase also introduces Audio Caption and Instruct TTS data. Relative to the earlier scaling stage, the cooldown phase emphasizes higher-quality multimodal supervision and longer-context capability refinement.

The technical consequence of this recipe is that the model learns more than a raw association between audio and text. It learns an operational interface between them. That interface is later reused in three directions: ASR maps audio evidence into text tokens, TTS maps text-side semantics into audio tokens, and Realtime couples listening, reasoning, and response generation under turn-level latency constraints. In this sense, pretraining is not merely background context for the rest of the report; it is the central mechanism that explains why all three specializations can share one backbone.

% \subsection{Why One Foundation Supports Three Different Speech Tasks}

% At first glance, ASR, TTS, and Realtime dialogue seem to ask for different competencies. ASR must compress a continuous signal into a discrete transcript, TTS must expand a textual description into a rich acoustic realization, and Realtime dialogue must do both while maintaining low latency, persona stability, and conversational appropriateness. The shared-foundation perspective explains why this tension is only partial.

% For ASR, strong language modeling is valuable because speech recognition errors are not purely acoustic. Homophones, punctuation, code-switching, named entities, and long-range consistency all require a model that can reason over discourse and lexical plausibility, not just frame-level acoustics. For TTS, the mirror image holds. High-quality synthesis is not purely acoustic either: naturalness, style consistency, role interpretation, and emotion control require a model that can understand instructions, context, and intent. Realtime dialogue goes one step further: it must infer not only what was said, but how it was said and what kind of response would be appropriate in the moment. In all three cases, the hard part lies at the boundary between signal and semantics. A single audio-language foundation is therefore not a compromise between unrelated tasks. It is a way of learning the shared boundary once, then specializing the direction and deployment regime of inference.
